\section{Results}
\subsection{An editable interface supports distinct predictive regimes}
Temporal memory improved the compact backbone on both datasets and the stronger TeRDy backbone on ICEWS05-15 (Table~\ref{tab:prediction}; Fig.~\ref{fig:prediction}a,b). For compact temporal factorization (CTF), mean filtered MRR increased from 0.3678 to 0.4311 on ICEWS14 and from 0.4546 to 0.5051 on ICEWS05-15. The paired gains were 6.32 percentage points (95\% interval 5.58--7.14) and 5.05 points (4.29--5.76), respectively. TeRDy achieved 0.6328 on ICEWS05-15, compared with 0.6285 for the neural anchor, a gain of 0.44 points (0.20--0.69). On ICEWS14, hybrid and neural MRR were 0.6616 and 0.6627, with a paired interval spanning zero. The interface therefore accommodates both substantial memory-driven gains and a predominantly neural predictive regime.

Validation-fitted belief models improved the calibration of the exposed hybrid answer (Fig.~\ref{fig:prediction}c--f). Across the four backbone--dataset combinations, mean Brier scores decreased from 0.188, 0.188, 0.223 and 0.240 for the raw hybrid maximum to 0.170, 0.182, 0.185 and 0.184, respectively. Mean ECE decreased to 0.025, 0.032, 0.027 and 0.022. This calibration supplies the correctness component of the acceptance policy; the certificate supplies its edit-stability component.

\begin{table*}[tbp]\centering\small\caption{Predictive performance on the fixed public test panels. Values are mean $\pm$ sample SD over three training seeds, in percent. MRR and Hits average over designated held-out answers; accuracy evaluates the selected object against the complete known-answer set. CTF denotes compact temporal factorization.}\label{tab:prediction}\begin{tabular}{lllrrrrr}\toprule Backbone & Dataset & Interface & MRR & Hits@1 & Hits@3 & Hits@10 & Accuracy\\\midrule
TeRDy & ICEWS14 & Neural & $66.27\pm0.14$ & $57.93\pm0.22$ & $71.49\pm0.26$ & $81.48\pm0.25$ & $58.28\pm0.20$\\
TeRDy & ICEWS14 & Hybrid & $66.16\pm0.10$ & $57.92\pm0.26$ & $71.17\pm0.09$ & $81.38\pm0.13$ & $58.29\pm0.24$\\
TeRDy & ICEWS05-15 & Neural & $62.85\pm0.70$ & $53.04\pm0.92$ & $68.96\pm0.63$ & $80.97\pm0.32$ & $53.47\pm1.00$\\
TeRDy & ICEWS05-15 & Hybrid & $63.28\pm0.92$ & $53.44\pm1.25$ & $69.33\pm0.82$ & $81.64\pm0.29$ & $53.85\pm1.24$\\
CTF & ICEWS14 & Neural & $36.78\pm0.42$ & $26.66\pm0.36$ & $41.83\pm0.81$ & $57.10\pm0.34$ & $27.05\pm0.30$\\
CTF & ICEWS14 & Hybrid & $43.11\pm0.20$ & $33.58\pm0.07$ & $48.01\pm0.43$ & $61.93\pm0.48$ & $34.05\pm0.10$\\
CTF & ICEWS05-15 & Neural & $45.46\pm0.08$ & $32.53\pm0.05$ & $51.98\pm0.51$ & $72.43\pm0.43$ & $32.98\pm0.03$\\
CTF & ICEWS05-15 & Hybrid & $50.51\pm0.22$ & $38.81\pm0.17$ & $56.64\pm0.41$ & $74.21\pm0.65$ & $39.28\pm0.11$\\\bottomrule\end{tabular}\end{table*}
\begin{figure*}[tbp]
\centering\includegraphics[width=\textwidth]{fig02_prediction.pdf}
\caption{\textbf{Predictive effects and answer calibration across two backbones.}
\textbf{a,} Neural and hybrid MRR on the fixed 5,000-query test panels. Points and bars show three-seed means and sample SD; connecting lines pair the two interfaces.
\textbf{b,} Paired hybrid-minus-neural MRR with 95\% timestamp-block bootstrap intervals from 2,000 draws, conditional on the three fitted models.
\textbf{c--f,} Reliability diagrams for the uncalibrated hybrid maximum and the validation-selected belief model. Each panel pools 15,000 query--seed predictions into ten equal-width probability bins; empty bins are omitted. The dashed diagonal denotes agreement between predicted probability and observed known-answer accuracy. CTF, compact temporal factorization.}
\label{fig:prediction}
\end{figure*}

\subsection{Certified selection exposes an explicit tolerance to deletion}
The exact one-window certificate retained a substantial pool of admissible predictions. Mean certified capacity was 96.75\% and 89.25\% for TeRDy on ICEWS14 and ICEWS05-15, and 59.16\% and 55.14\% for CTF. Some queries had no admissible history: 28.88\% on ICEWS14 and 13.30\% on ICEWS05-15. Restricting the calculation to queries with evidence still gave certified fractions of 95.43\%, 87.60\%, 42.58\% and 48.26\%, respectively. Certification therefore extended beyond the empty-memory fallback.

Applying the certificate as a gate eliminated allowed worst-window flips among accepted answers (Fig.~\ref{fig:selection}; Table~\ref{tab:selection}). Across 60,000 query--seed cases and three deletion budgets, all 180,000 direct witness checks agreed with the certificate, with zero certified violations. At 40\% coverage, the CTF belief selector exposed answers with worst-window flip rates of 32.15\% on ICEWS14 and 10.70\% on ICEWS05-15; the certified policy reduced both rates to zero. Its answer-error rates were 47.78\% and 36.57\%, compared with 46.43\% and 35.28\% for belief alone. The joint rate of an incorrect initial answer or a worst-window flip decreased by 11.07 points (95\% interval 9.36--12.89) and 4.05 points (3.59--4.43). These results quantify the additional stability obtained at the selected coverage and its accompanying correctness trade-off.

Confidence selection already retained almost entirely stable TeRDy answers at low coverage. Its value from the certificate became more visible as coverage increased: at 80\% coverage, joint error decreased by 0.28 points on ICEWS14 and 0.96 points on ICEWS05-15. For the latter interval, 1,748 of 2,000 bootstrap draws retained enough certified queries in every seed, as detailed in the supplement. At the same time, certificate-feature calibration alone had higher mean AURC than the strongest validation-selected control in all four combinations. The gate consequently serves as a defined stability requirement, with confidence and ensemble controls retained as the appropriate comparators for correctness ranking.

\begin{table*}[tbp]\centering\small\caption{Calibrated belief and one-window certified selection on the same predictions. Error and edit-flip rates are percentages at exactly 20\% or 40\% query coverage. Values are three-seed mean $\pm$ SD. Certified capacity is the fraction of all queries passing the certificate. Unavailable coverage is denoted by a dash.}\label{tab:selection}\begin{tabular}{lllrrrr}\toprule Backbone & Dataset & Policy & Error at 20\% & Error at 40\% & Flips at 40\% & Capacity\\\midrule
TeRDy & ICEWS14 & Belief & $8.93\pm3.27$ & $14.70\pm2.70$ & $0.00\pm0.00$ & --\\
TeRDy & ICEWS14 & Certified & $8.93\pm3.27$ & $14.70\pm2.70$ & $0.00\pm0.00$ & $96.75\pm2.08$\\
TeRDy & ICEWS05-15 & Belief & $10.07\pm1.22$ & $18.70\pm0.54$ & $0.03\pm0.06$ & --\\
TeRDy & ICEWS05-15 & Certified & $10.07\pm1.22$ & $18.68\pm0.52$ & $0.00\pm0.00$ & $89.25\pm7.22$\\
CTF & ICEWS14 & Belief & $30.90\pm0.61$ & $46.43\pm0.42$ & $32.15\pm1.66$ & --\\
CTF & ICEWS14 & Certified & $30.97\pm0.81$ & $47.78\pm0.59$ & $0.00\pm0.00$ & $59.16\pm1.64$\\
CTF & ICEWS05-15 & Belief & $21.43\pm0.83$ & $35.28\pm0.51$ & $10.70\pm3.71$ & --\\
CTF & ICEWS05-15 & Certified & $21.43\pm0.83$ & $36.57\pm1.05$ & $0.00\pm0.00$ & $55.14\pm2.18$\\\bottomrule\end{tabular}\end{table*}
\begin{figure*}[tbp]
\centering\includegraphics[width=\textwidth]{fig03_selection.pdf}
\caption{\textbf{Correctness and resistance to evidence deletion are separate selection objectives.}
\textbf{a--d,} Answer-error risk--coverage curves for the best validation-selected control, the one-encoder calibrated belief policy and its one-window certified gate.
\textbf{e--h,} The corresponding worst-case one-window flip rates, evaluated on exactly the same selected predictions. Curves show three-seed means; bands show sample SD. A certified curve stops at the largest plotted coverage reachable by every seed. Controls include temperature scaling and a three-seed ensemble, and their identities are frozen on validation data. Each panel uses all 5,000 test queries per seed.}
\label{fig:selection}
\end{figure*}

\subsection{Certificate surplus predicts vulnerability beyond the original partition}
The surplus retained diagnostic value when the intervention changed (Fig.~\ref{fig:diagnostics}). For random two-window deletions, shifted window boundaries, doubled window width and single-event deletion, its AUROC exceeded that of the original hybrid margin in every backbone--dataset combination. The 16 paired gains ranged from 0.064 to 0.584, and every 95\% interval lay above zero. On TeRDy, where margin was already informative, the surplus reached AUROCs of 0.991--0.998 on ICEWS14 and 0.970--0.992 on ICEWS05-15. For CTF, the corresponding ranges were 0.903--0.969 and 0.847--0.982.

This diagnostic comparison concerns direct interventions on the same fixed predictions. A large current score gap can arise from highly concentrated memory support, whereas the surplus accounts for how that gap changes when shared evidence is removed. The conceptual equal-score example in Fig.~\ref{fig:interface} isolates this distinction exactly, and the public-data perturbations show its empirical relevance across different evidence partitions. The formal guarantee remains tied to the declared window family; the alternative interventions measure how informative the score remains outside that exact family.

\begin{figure*}[tbp]
\centering\includegraphics[width=\textwidth]{fig04_diagnostics.pdf}
\caption{\textbf{The certificate surplus identifies edit-sensitive answers more accurately than the original score margin.}
\textbf{a--d,} AUROC differences between negative one-window surplus and negative hybrid margin for the four indicated backbone--dataset combinations. The outcomes are an answer change under at least one of eight random two-window deletions, a worst-case shifted seven-unit window deletion, a worst-case 14-unit window deletion, and removal of the latest winner-supporting event. Each comparison pools 15,000 query--seed predictions. Bars show paired 95\% intervals from 2,000 timestamp-block bootstrap draws, resampled jointly across the three fixed fits.}
\label{fig:diagnostics}
\end{figure*}

\subsection{Exact certification resolves the cost of conservative and sampled checks}
Exact certification recovered admissible answers that a conservative bound rejected (Fig.~\ref{fig:controls}a). Among queries with nonempty memory, its one-window certified coverage exceeded the winner-only bound by 2.45 and 12.09 percentage points for TeRDy on ICEWS14 and ICEWS05-15, and by 2.28 and 2.49 points for CTF. The largest difference arose on TeRDy/ICEWS05-15, where coverage increased from 80.39\% to 92.48\%. Accounting for the competitor's simultaneous evidence loss therefore makes the guarantee more useful without relaxing the deletion budget.

Finite deletion probes had a different limitation (Fig.~\ref{fig:controls}b). On CTF/ICEWS05-15, 31.93\% of nonempty-memory queries that passed eight probes were vulnerable to another allowed deletion. Increasing the probe count to 32 and 64 reduced this fraction to 18.60\% and 11.35\%, respectively. At 64 probes, the corresponding fractions were 0\%, 0.67\% and 1.57\% for TeRDy/ICEWS14, TeRDy/ICEWS05-15 and CTF/ICEWS14. The exact certificate resolves this residual uncertainty by covering every allowed subset analytically.

The prespecified weight sweep explains why the predictive operating point is selected separately from the deletion budget (Fig.~\ref{fig:controls}c,d). Greater memory weight increased CTF MRR over its neural-only value throughout the grid, while the strongest TeRDy performance occurred at smaller weights. Certified capacity decreased as memory acquired more influence in these measurements. The complete sweep, including $\lambda=0$, is reported in the supplement. A fixed-denominator calculation disagreed with the renormalized one-window certificate on 0.14--8.22\% of all queries, demonstrating that the edited-score definition also matters.

Sparse competitor reduction made the exact calculation practical. Over the 32-query microbenchmark, including empty-memory queries, median paired dense-to-sparse runtime ratios were 4.63--11.75 across the four combinations. On nonempty histories they were 7.50--17.52. These timings compare a sparse implementation returning certificates, radius and witnesses with a dense implementation returning certificate indicators; all-query and nonempty-history timings are reported separately in the supplement. Together, the controls distinguish exact certification from both conservative abstention and finite intervention sampling.

\begin{figure*}[tbp]
\centering
\includegraphics[width=\textwidth]{fig05_controls.pdf}
\caption{\textbf{Analytical exactness and the evidence-weight trade-off.}
\textbf{a,} Exact and conservative one-window certificate coverage among queries with nonempty memory.
\textbf{b,} The fraction of predictions passing all sampled deletion probes that are vulnerable to the exact one-window adversary, for 8, 32 and 64 probes.
\textbf{c,d,} Filtered MRR and one-window certified coverage over the prespecified evidence-weight grid. These controls use seed 0 of every backbone--dataset combination and all 5,000 test queries; temperature is held at its validation-selected value.}
\label{fig:controls}
\end{figure*}

\subsection{Local maintenance restricts recomputation to affected answers}
Local maintenance reduced the number of refreshed query states from 1,200,000 to 55,064 across all 240 transactions, a reduction of 95.41\% (Fig.~\ref{fig:maintenance}). The reductions were 95.59\%, 94.21\% and 96.44\% on ICEWS14, ICEWS05-15 and GDELT, respectively. Both implementations applied identical event writes and maintained identical window statistics; the difference was whether every registered query or only the affected set was refreshed. Their outputs agreed after every transaction, and independent raw-event rebuilds agreed on all discrete outputs for 128 queries per dataset, with maximum absolute numerical difference $2.22\times10^{-16}$.

The benefit remained measurable at the largest transaction size. For 1,000 edited events, median local refresh times were 0.330, 2.407 and 0.504 seconds on the three datasets, compared with 1.367, 4.030 and 2.161 seconds for full refresh. The median paired speedups were 4.18, 1.67 and 4.29, respectively. Median affected sets contained 739, 973.5 and 631.5 of the 5,000 registered queries. Smaller transactions touched fewer answers and produced larger timing ratios. Uniformly sampled single-event edits affected no registered query in 80\%, 60\% and 20\% of transactions, respectively; the supplement reports these frequencies alongside every timing stratum.

Initial cache construction took 2.20--5.46 seconds. Numeric arrays occupied 283, 700 and 275~MiB on ICEWS14, ICEWS05-15 and GDELT; measured process-memory increments, including container overhead, were 434, 951 and 305~MiB. These one-time costs are separated from transaction timings. The results show that explicit evidence dependencies support both the stability calculation and its efficient maintenance, including on the larger GDELT event history.

\begin{figure*}[tbp]
\centering
\includegraphics[width=\textwidth]{fig06_maintenance.pdf}
\caption{\textbf{Maintaining only affected answers reduces update work.}
\textbf{a--c,} Full and local refresh times for 5,000 registered queries after transactions of 1, 10, 100 or 1,000 events. Points show medians and bands show interquartile ranges over ten sampled deletion/restoration pairs at each batch size, with alternating implementation order.
\textbf{d--f,} The corresponding fraction of registered queries refreshed. The full control refreshes 100\% at every transaction. Timings include event writes and certificate maintenance; encoder inference and initial cache construction are excluded. GDELT uses a fixed training relation-frequency anchor and tests maintenance scaling.}
\label{fig:maintenance}
\end{figure*}

